\documentclass[11pt]{article}
\usepackage[hmargin=0.85in,vmargin=0.8in]{geometry}
\usepackage{amsmath,amssymb,amsthm}
\usepackage[T1]{fontenc}
\usepackage{lmodern}
\usepackage[hidelinks]{hyperref}
\setlength{\emergencystretch}{1em}
\newtheorem{theorem}{Theorem}
\newtheorem{lemma}{Lemma}
\newcommand{\C}{\mathbb C}
\newcommand{\R}{\mathbb R}
\newcommand{\Tr}{\operatorname{Tr}}
\title{Three Coordinate Quasitraces Refute Conjecture 00000000918}
\author{idealistichacker \quad (AI-assisted)}
\date{October 30, 2026}
\begin{document}
\maketitle
\begin{abstract}
The printed assertion that the quasitrace simplex of every nuclear algebra
has dimension one fails already for the finite-dimensional commutative
C*-algebra $A=\C^3$ with its supremum norm. Its three coordinate evaluations
are positive normalized weakly continuous traces with genuine compatible
matrix quasitrace extensions. They are affinely independent over $\R$.
The full quasitrace space containing them therefore cannot be one-dimensional.
We formalize an exact completely positive approximation witness, the
quasitrace and matrix-extension laws, the full-space dictionary, and the
closed negation of the necessary universal dimension clause.
\end{abstract}

\section{The Clause Being Refuted}
Both source languages assert a simplex structure for the quasitrace space
of a nuclear algebra, identify its extreme points with factor traces, and
require that the same simplex have dimension one and a closed extreme-point
set. Neither imposes simplicity, factoriality of the algebra, or restriction
to a selected one-parameter family. We refute only the universal
\emph{dimension-one} requirement. We do not assert that the other clauses
are false.

We use the usual nuclear C*-algebra and complex-valued quasitrace setting.
In particular, dimension of a trace simplex means the real dimension of
its affine span, not the dimension of an individual extreme-point factor.
The usual normalized convention is $\tau(1)=1$. Omitting normalization
does not remove the three witnesses below.

\begin{theorem}\label{thm:main}
There exists a commutative, finite-dimensional, unital nuclear C*-algebra
whose full normalized weakly continuous $2$-quasitrace space has real affine
dimension at least two. The corresponding space without normalization also
does not have affine dimension one.
\end{theorem}
Since this is a genuine nuclear algebra, a universal assertion for nuclear
algebras must apply to it. Restricting the printed clause to finite-dimensional
commutative nuclear algebras is already enough to obtain a false statement.

\section{The Actual Algebra and Nuclearity Witness}
Let $A$ be the functions on $\{0,1,2\}$ with values in $\C$, with pointwise
operations, coordinate conjugation, unit $(1,1,1)$, and
\[
 \|a\|=\max_{0\leq i<3}|a_i|.
\]
This is the ordinary supremum-norm C*-algebra $\C^3$, not Euclidean space
with its $\ell^2$ norm. Its complex dimension is three.

Use the standard completely positive approximation property (CPAP)
characterization of nuclearity: for every finite $E\subset A$ and
$\varepsilon>0$, there is a finite-dimensional C*-algebra $F$ and completely
positive contractions $\phi:A\to F$, $\psi:F\to A$ such that
\[
 \|\psi(\phi(a))-a\|<\varepsilon\qquad(a\in E).
\]
Take $F=A$ and $\phi=\psi=\mathrm{id}_A$. Every matrix amplification of
the identity is the actual identity, hence preserves positivity; both maps
are contractions and every error is zero. Thus the full finite-set and
positive-error quantifiers hold, rather than being assumed from a label
attached to $A$. The standard CPAP/nuclearity dictionary is the mathematical
convention used here \cite{CE,CT}. The Lean result proves CPAP directly; it
does not claim a new formal proof of its tensor-norm equivalence.

\section{Genuine Quasitraces and Matrix Extensions}
For $i\in\{0,1,2\}$ put $\tau_i(a)=a_i$. This is an actual continuous
complex-linear functional, with $\tau_i(1)=1$ and $|\tau_i(a)|\leq\|a\|$.
Moreover,
\[
 \tau_i(ab)=\tau_i(ba),\qquad
 \tau_i(a^*a)=\tau_i(aa^*)=|a_i|^2\geq0,\qquad
 \tau_i(a^*)=\overline{\tau_i(a)}.
\]
Thus each is a positive normalized trace.

For clarity, a standard $1$-quasitrace has positive cyclic star-square
values, real values on selfadjoint elements, compatibility with the
selfadjoint real/imaginary decomposition, and linearity on abelian
C*-subalgebras. A $2$-quasitrace additionally has a compatible $1$-quasitrace
on $M_2(A)$ \cite{H}. Every positive complex-linear star-preserving trace
satisfies these laws. In particular, additivity gives the source's
subadditivity with equality on the selfadjoint/positive part, where the
order inequality is defined.

Here the compatible extension is explicit. For $X\in M_2(A)$ set
\[
 Q_i(X)=\Tr_2\bigl[(\tau_i(X_{rs}))_{r,s=0,1}\bigr]
       =\tau_i(X_{00})+\tau_i(X_{11}).
\]
This uses the \emph{unnormalized} ordinary matrix trace. It is complex-linear,
star-preserving and tracial, and
\[
 Q_i(X^*X)=\sum_{r,s=0}^1 |(X_{rs})_i|^2\geq0.
\]
It therefore satisfies the full $1$-quasitrace laws on the actual matrix
C*-algebra. For the standard corner embedding,
\[
 Q_i\begin{pmatrix}a&0\\0&0\end{pmatrix}=\tau_i(a),
 \qquad Q_i(1_2)=2.
\]
The value two on the matrix identity is correct; imposing normalization one
on this extension would break corner compatibility. Thus the $\tau_i$
are genuine normalized $2$-quasitraces, not three functions merely named
quasitraces.

For $\sigma(A,A^*)$, every continuous linear functional is weakly
continuous by the definition of the weak topology. Lean checks this
on its actual \texttt{WeakSpace} for each $\tau_i$. In finite dimension,
the ordinary weak and norm topology descriptions do not reduce these
three distinct functionals to a one-parameter family. The global theorem
that quasitraces on exact algebras are traces is not needed.

\section{The Full Space and Its Affine Dimension}
Let $p_j$ be the coordinate projection with value one at $j$ and zero
elsewhere. These are actual elements of $A$, and $\tau_i(p_j)=\delta_{ij}$.
If
\[
 r(\tau_1-\tau_0)+s(\tau_2-\tau_0)=0\qquad(r,s\in\R),
\]
evaluation at $p_1$ and $p_2$ gives $r=0$ and $s=0$. Hence the three
coordinate traces are affinely independent over $\R$.

This is an obstruction for the \emph{entire} space. Let $S$ be any set
in the real continuous dual containing all three $\tau_i$. Its affine-span
direction contains the two independent differences, so its dimension is
at least two. Taking $S$ to be the full normalized weak $2$-quasitrace
space proves Theorem~\ref{thm:main}. Taking the full space without
normalization proves the second assertion. The set $S$ is not defined to
be the three displayed points or their convex hull.

There is no hidden restriction to preselected linear functionals: on this
commutative algebra, the quasitrace axiom of linearity on the whole closed
abelian subalgebra $A$ gives an actual complex-linear functional. Finite
dimension makes it continuous. Lean proves its unique representation in
$A\to_L\C$, and identifies the full unbundled normalized space with the
image of that set under an injective real-linear coefficient map. Thus
the dual is the natural affine ambient space for the original functionals.
Its real finite-dimensionality is established before using \texttt{finrank};
no infinite-dimensional \texttt{finrank=0} convention is exploited.

We only require dimension at least two, not an exact dimension computation,
an affine-homeomorphism classification, or a theorem about all nuclear
algebras. The factor-trace and closed-extreme-set clauses are not attacked.

\section{Formal Scope and Evidence Boundaries}
The five Lean modules use Lean 4.33.1 and fixed Mathlib revision
\texttt{0df444a360eaa60ab8c11dca51a86af692955474}. They separate the exact
CPAP witness, genuine Pi-norm algebra and continuous-dual geometry, actual
operator-norm \texttt{CStarMatrix} extension, quasitrace axioms and full-space
dictionary, and final closed negation
\texttt{source\_necessary\_clause\_false}. The matrix base order is the
canonical spectral order of the same C*-algebra, not a replacement of its
norm by the raw entrywise matrix norm.

All 68 public theorem types and transitive axiom footprints have actual
author-side compiler output. Allowed logical axioms are only
\texttt{propext}, \texttt{Classical.choice}, and \texttt{Quot.sound}.
Failed compilation attempts, including an initial 60-second CPAP timeout,
are retained separately and are not counted as successes. A small exact
rational-complex diagnostic checks coordinate arithmetic and a nonzero
affine minor; those finite checks do not establish nuclearity, general
quasitrace laws, weak continuity or the all-input Lean statements.

Original source blob:
\texttt{dadc2ef65f6a57feb3def12febafd66299b1f918}, checked in both languages.
This research manuscript and author-side module checks do not constitute
independent frozen-package/PDF review, portable whole-package reproduction,
current duplicate clearance or upstream acceptance. Existing compiler and
official dependency caches remain trust inputs; no full-dependency-clean,
empty-machine, Linux or hosted-CI result is claimed.

\begin{thebibliography}{9}
\raggedright
\bibitem{H} U. Haagerup,
\emph{Quasitraces on exact C*-algebras are traces},
C. R. Math. Rep. Acad. Sci. Canada 36 (2--3) (2014), 67--92;
arXiv:1403.7653. Introduction: quasitrace and compatible matrix-extension
conventions. The global theorem is not used as an assumption.
\bibitem{CT} J. Castillejos, S. Evington, A. Tikuisis, S. White and W. Winter,
\emph{Nuclear dimension of simple C*-algebras}, Inventiones Mathematicae
(2021), DOI 10.1007/s00222-020-01013-1, Section 1.1.
\bibitem{CE} M.-D. Choi and E. G. Effros,
\emph{Nuclear C*-algebras and the approximation property},
American Journal of Mathematics 100 (1) (1978), 61--79.
Foundational bibliographic reference, not a claim to have independently
audited the entire original equivalence proof.
\end{thebibliography}
\end{document}
